\section*{Chapter \ref{ch:g12:fiscal-policy-and-debt} --- Fiscal Policy and Public Debt}
% figdata: figdata/grade-12/11-fiscal-policy-and-debt/out/ (debt_long, debt_euro, uk_decomp, debt_key).

\begin{solution}{exo:g12:fiscal-policy-and-debt:1}
$-4 + 3 = -1$: a primary deficit of 1\,\% of GDP.
\end{solution}

\begin{solution}{exo:g12:fiscal-policy-and-debt:2}
\omterm{def:g12:fiscal-policy-and-debt:primary-balance}{Primary balance} $40 - 38 = +2$; overall balance $2 - 5 = -3$\,\% of GDP.
\end{solution}

\begin{solution}{exo:g12:fiscal-policy-and-debt:3}
In 1946, at 269.8\,\%; 46.7\,\% in 1975.
\end{solution}

\begin{solution}{exo:g12:fiscal-policy-and-debt:4}
118.9\,\% in 1946 and 128.1\,\% in 2020: higher in 2020 than at the end of the Second World War.
\end{solution}

\begin{solution}{exo:g12:fiscal-policy-and-debt:5}
$60 \times 1.04/1.03 \approx 60.6\,\%$.
\end{solution}

\begin{solution}{exo:g12:fiscal-policy-and-debt:6}
$a = 1.02/1.05 \approx 0.9714$: $d_1 \approx 97.14 + 1 = 98.14$, $d_2 \approx 95.34 + 1 = 96.34$, $d_3 \approx 93.59 + 1 = 94.59\,\%$.
\end{solution}

\begin{solution}{exo:g12:fiscal-policy-and-debt:7}
$a = 1.06/1.03 \approx 1.0291$: 102.9, 105.9 and 109.0\,\%.
\end{solution}

\begin{solution}{exo:g12:fiscal-policy-and-debt:8}
$120 \times (0.05 - 0.03)/1.03 \approx 2.33\,\%$ of GDP, a primary surplus.
\end{solution}

\begin{solution}{exo:g12:fiscal-policy-and-debt:9}
From $100 \times (-0.01)/1.04 \approx -0.96$ (a primary deficit is allowed) to $100 \times 0.01/1.04 \approx +0.96$ (a surplus is needed): a change of about 1.9 points of GDP.
\end{solution}

\begin{solution}{exo:g12:fiscal-policy-and-debt:10}
Spain: from 35.7 to 100.5\,\%, a rise of 64.8 points (France: 50.2 points). Ireland's fell most after 2012, from 119.0 to 34.8\,\%, by 84.2 points; Germany's also fell, from 79.8 to 62.8\,\% ($-17.0$ points).
\end{solution}

\begin{solution}{exo:g12:fiscal-policy-and-debt:11}
Receipts fall by $0.40 \times 3 = 1.2$\,\% of GDP, benefits rise by 0.3: the deficit widens by about 1.5 points.
\end{solution}

\begin{solution}{exo:g12:fiscal-policy-and-debt:12}
Households that received the 2008 rebates spent about 50--90\,\% of them within three months, counting durables, instead of saving them for future taxes. It points to the
assumption that households can borrow freely and look far ahead: many spend extra cash when it arrives.
\end{solution}

\begin{solution}{exo:g12:fiscal-policy-and-debt:13}
$d_{t+1} - d^* = (a d_t - p) - (a d^* - p) = a (d_t - d^*)$, so by induction $d_t - d^* = a^t (d_0 - d^*)$. With $0 < a < 1$, $a^t \to 0$ and $d_t \to d^*$; with $a > 1$ the gap
grows without limit unless $d_0 = d^*$.
\end{solution}

\begin{solution}{exo:g12:fiscal-policy-and-debt:14}
\omterm{def:g10:national-accounts:nominal-gdp}{Nominal GDP} rises by 10\,\% while the debt does not: $100/1.1 \approx 90.9\,\%$. Bondholders lose: their \omterm{def:g12:present-value:bond}{bonds} buy about 9\,\% less.
\end{solution}

\begin{solution}{exo:g12:fiscal-policy-and-debt:15}
The private half, 85\,\% of GDP, is cut by $0.535 \times 85 \approx 45.5$ points: the ratio falls to about 124.5\,\%. It may rise again if output keeps falling, if deficits
continue, or if the remaining debt pays high rates.
\end{solution}

\begin{solution}{pb:g12:fiscal-policy-and-debt:1}
\textbf{1.} $269.8 - 46.7 = 223.1$ points.

\textbf{2.} $1.0305 \times 1.0582 \approx 1.090$: about 9.0\,\%, as observed.

\textbf{3.} $270 \times 1.03/1.09 \approx 255.1\,\%$.

\textbf{4.} $p = 270 \times (0.03 - 0.09)/1.09 \approx -14.9$: a primary deficit of 14.9\,\% of GDP would have kept the ratio constant.

\textbf{5.} Below: 6.42\,\% against 9.0\,\%. The ratio falls even with a balanced primary budget.

\textbf{6.} 189.5 points.

\textbf{7.} $104.8 + 192.2 = 297.0$ points.

\textbf{8.} $297.0 - 189.5 = 107.5$ points.

\textbf{9.} $107.5/223.1 \approx 48\,\%$.

\textbf{10.} Primary surpluses and other changes in the debt (sales of assets, revaluations, differences between the data and the approximation): 115.7 points, 52\,\%.

\textbf{11.} Prices raise \omterm{def:g10:national-accounts:nominal-gdp}{nominal GDP}, while the debt, fixed in money, does not rise with them.

\textbf{12.} The holders of the \omterm{def:g12:present-value:bond}{bonds}, whose interest and repayment bought less.

\textbf{13.} The debt consists of \omterm{def:g12:present-value:bond}{bonds} issued in earlier years at the rates of those years, and of short-term debt at other rates; the average rate paid differs from
the current long rate.

\textbf{14.} As inflation rose, lenders expected more of it and asked for compensation in the rate.

\textbf{15.} $118.9 - 33.6 = 85.3$ points.

\textbf{16.} Real growth, inflation and primary surpluses; not default.

\textbf{17.} Default: private bondholders accepted a 53.5\,\% cut in face value.

\textbf{18.} Yes: with $a < 1$ and a constant primary deficit, the ratio tends to $d^* = -p/(1 - a)$, a finite level.

\textbf{19.} Inflation removed 192 of the 223 points; with low inflation nominal growth would have been lower, and interest rates too perhaps; the record shows one path
only.

\textbf{20.} Nominal growth above the interest rate accounts for 107.5 of the 223.1 points, about 48\,\% of the fall; the rest came from primary surpluses and from the gap between the long-term rate and the rate actually paid.
\end{solution}
